\PassOptionsToPackage{table,xcdraw}{xcolor}
\documentclass[10pt, a4paper]{article}

\usepackage[T1]{fontenc}
\usepackage[utf8]{inputenc}
\usepackage{tgpagella}
\usepackage{tgcursor}
\usepackage[spanish, es-noshorthands]{babel}
\usepackage{amsmath, amssymb}
\usepackage{booktabs}
\usepackage{tabularx}
\usepackage{xcolor}
\usepackage{tikz}
\usetikzlibrary{shapes.geometric, arrows.meta, positioning}
\usepackage[american]{circuitikz}
\usepackage[most]{tcolorbox}
\usepackage[margin=1.5cm]{geometry}
\usepackage{fancyhdr}

\setlength{\headheight}{16pt}
\setlength{\emergencystretch}{2em}
\definecolor{ucbblue}{RGB}{0, 51, 102}
\definecolor{neutralbg}{RGB}{247, 248, 250}
\definecolor{warnbg}{RGB}{254, 246, 246}
\newcolumntype{P}[1]{>{\raggedright\arraybackslash}p{#1}}
\newcolumntype{Y}{>{\raggedright\arraybackslash}X}

\pagestyle{fancy}
\fancyhf{}
\lhead{\textbf{UCB Sede Tarija} | IMT-221}
\chead{Lab 5: Rectificación de Precisión (Rutas A y B)}
\rhead{Semana 14 - Sesión 2}
\renewcommand{\headrulewidth}{0.4pt}

\begin{document}

\begin{center}
    {\Large\textbf{\textcolor{ucbblue}{Guía de Trabajo y Registro: Laboratorio 5 (Rectificación)}}} \\[0.1cm]
    {\footnotesize\textbf{Objetivos:} 1. Diodo real/señal pequeña; 2. Superdiodo y estados; 3. Predecir límites; 4. Evidencia verificable.}
\end{center}

\section*{1. Modelado del Diodo Convencional y Superdiodo No Inversor (V1 y V3)}
\begin{center}
    \begin{circuitikz}[>=Stealth, thick, scale=0.85, transform shape]
        % Diodo Convencional
        \node[font=\bfseries, text=ucbblue] at (1.5, 3.2) {Convencional};
        \draw (0,0) node[ground]{} to[sV, v={$v_i(t)$}] (0,2) to[short] (1,2) to[D, l=$D$, i=$i_D$] (3,2) to[short, -o] (4,2) node[right]{$v_{o,D}$};
        \draw (3,2) to[R, l=$R_L$] (3,0) node[ground]{};
        
        % Superdiodo No Inversor
        \node[font=\bfseries, text=ucbblue] at (9, 3.2) {De Precisión (Superdiodo)};
        \draw (7,0) node[op amp] (op) {};
        \draw (op.+) to[short, -o] ++(-1,0) node[left]{$v_i(t)$};
        \draw (op.out) to[short, *-] ++(0.5,0) coordinate(va) node[above]{$v_a$} to[D, l=$D$] ++(2,0) coordinate(vo) to[short, *-o] ++(0.5,0) node[right]{$v_{o,P}$};
        \draw (vo) to[R, l=$R_L$] ++(0,-1.5) node[ground]{};
        \draw (op.-) -- ++(0,1) -| (vo);
        \draw (op.up) -- ++(0,0.5) node[above]{\footnotesize $V_{S+}$};
        \draw (op.down) -- ++(0,-0.5) node[below]{\footnotesize $V_{S-}$};
    \end{circuitikz}
\end{center}

\textbf{Modelo de umbral (Convencional):} $v_o(t) \approx \max\{v_i(t) - V_\gamma, 0\}$. Falla con amplitudes menores o cercanas a $V_\gamma$. \\
\textbf{Estados del Superdiodo (Precisión):}
\begin{table}[htbp]
    \centering
    \footnotesize
    \renewcommand{\arraystretch}{1.3}
    \begin{tabularx}{\textwidth}{|P{2.2cm}|c|P{4.2cm}|X|X|}
    \hline
    \rowcolor[HTML]{EFEFEF} 
    \textbf{Condición} & \textbf{Estado Diodo} & \textbf{Estado del Lazo} & \textbf{Salida Rectificada ($v_{o,P}$)} & \textbf{Salida Interna ($v_a$)} \\ \hline
    $v_i > 0$ (lineal) & Conduce & Cerrado ($v_- \approx v_+ = v_i$) & $v_o \approx v_i$ (Copia fiel) & $v_a \approx v_o + v_D$ \\ \hline
    $v_i < 0$ & Bloquea & \textbf{Abierto} ($v_- \neq v_+$) & $v_o \approx 0$ (por $R_L$) & Tiende al límite negativo. \\ \hline
    \end{tabularx}
\end{table}

\textbf{E1 / E3 (Diagnóstico Rápido):} Un diodo convencional con $V_\gamma = 0.40\,\text{V}$ ante un estímulo de $0.30\,\text{V}$ pico predice $0\,\text{V}$ de salida (no supera umbral). El superdiodo soluciona esto porque, al conducir, la realimentación fuerza a $v_a$ a sobreponerse a la caída del diodo. No asuma $v_- = v_+$ cuando el diodo bloquea (lazo abierto).

\section*{2. Decisión Operativa (V5) y Cálculo de Error Experimental}
\begin{center}
    \begin{tikzpicture}[>=Stealth, thick, scale=0.85, transform shape]
        \node[draw=ucbblue, rectangle, fill=ucbblue!8, text width=6.5cm, align=center, minimum height=0.9cm] (q) at (0,0) {\textbf{¿Esquema, modelo de CI, fuente y manuales verificados?}};
        \node[draw=green!60!black, rectangle, fill=green!10, text width=4cm, align=center, minimum height=0.8cm] (rutaA) at (-3.8, -1.8) {\textbf{Ruta A: MEDIDO}\\Autorización a banco físico};
        \node[draw=orange, rectangle, fill=orange!10, text width=5cm, align=center, minimum height=0.8cm] (rutaB) at (3.8, -1.8) {\textbf{Ruta B: SIMULADO / PREDICHO}\\La simulación no es laboratorio físico.};
        \draw[->] (q) -- node[left, font=\scriptsize]{\textbf{SÍ}} (rutaA);
        \draw[->] (q) -- node[right, font=\scriptsize]{\textbf{NO} (Faltan datos)} (rutaB);
    \end{tikzpicture}
\end{center}
Para el cálculo de error (E6), la referencia ideal es $v_{\mathrm{ref}}(t) = \max\{v_i(t), 0\}$. El error punto a punto es $e(t) = v_o(t) - v_{\mathrm{ref}}(t)$. Defina $E_1 = \frac{1}{T_{\mathrm{obs}}} \int |e(t)|\, dt \quad [\mathrm{V}]$.

\section*{3. Hoja de Registro (Completar con Datos Reales o Simulación Verificada)}
\begin{table}[htbp]
    \centering
    \scriptsize
    \renewcommand{\arraystretch}{1.45}
    \begin{tabularx}{\textwidth}{|P{4.2cm}|Y|}
    \hline
    \rowcolor[HTML]{EFEFEF} 
    \textbf{Campo / Parámetro del Ensayo} & \textbf{Valor a Documentar, Archivo o Captura (Indicar unidades)} \\ \hline
    Grupo / Versión y Autores & \textbf{Evidencia:} $\square$ MEDIDO \quad $\square$ SIMULADO \quad $\square$ PREDICHO \quad $\square$ NO OBTENIDO \\ \hline
    Circuito / Op-Amp / Diodo / $R_L$ & \\ \hline
    Alimentación $V_{S+}, V_{S-}$ y Tierra Instr. & \\ \hline
    Fuente ($v_i$): Amplitud, $f$, offset & \\ \hline
    Resultado Convencional ($v_{o,D}$) & \\ \hline
    Resultado Precisión ($v_{o,P}$) y $v_a$ & \\ \hline
    Predicción vs. Observación (Ej. $E_1$) & \\ \hline
    Problema, Hipótesis y Próximo Test & \\ \hline
    \end{tabularx}
\end{table}
\noindent\textbf{Preguntas de Reflexión (E9):} a) ¿Por qué el superdiodo rectifica señales $\ll 0.7\,\text{V}$? (Por la realimentación). b) ¿Limitación de implementación? (SR, GBW y saturación/recuperación). c) ¿Evidencia de Lab 6 no obtenida aquí? (Bode espectral)[cite: 12].

\end{document}
